% §3.4 Miller-Rabin
%
% Cross-reference resolution status:
%
% RESOLVED:
%   \Cref{sec:and-under-noise}, \Cref{sec:composition-theorem},
%   \Cref{eq:composition-theorem}  -- chapter 3
%   \Cref{sec:axiom-1}, \Cref{ax:element-indep}, \Cref{def:bernoulli-bool} -- chapter 4
%   \Cref{sec:bridge-to-sets}      -- chapter 3 §3.5
%   \Cref{ch:randomized-algorithms}-- chapter 11
%
% DEFINED here: \label{sec:miller-rabin}, \label{eq:mr-fpr},
%   \label{eq:mr-k-round}, \label{alg:miller-rabin}
%
% Monograph register: Miller-Rabin as a k-fold one-sided Bernoulli[Bool]
% conjunction. PRESERVE the B1 fix: the k-round FPR is the product
% prod(eps_i) <= (1/4)^k, with the equality stated as the all-negative dual
% of the composition theorem.

\section{Miller-Rabin}
\label{sec:miller-rabin}

Miller-Rabin primality testing is a $k$-fold one-sided $\Bernoulli{\BoolType}$
conjunction. Number theory contributes one fact, the per-round bound $\fprate
\le 1/4$; everything else is the composition machinery of \S\S3.1--3.3.

\subsection*{The witness test as a noisy bit}

Fix an odd $n > 3$ and write $n - 1 = 2^r d$ with $d$ odd. A witness $a \in
\{2, \ldots, n-2\}$ certifies compositeness if the sequence $a^d, a^{2d},
\ldots, a^{2^r d} \pmod{n}$ reaches $1$ without passing through $-1$ and
without starting at $\pm 1$; otherwise the test reports ``probably prime.''
Identify ``probably prime'' with $\top$ and ``composite'' with $\bot$, and let
the true value be whether $n$ is prime. Then a prime always passes (by Fermat's
little theorem and the fact that the only square roots of $1$ modulo a prime
are $\pm 1$), so the test has no false negatives, $\fnrate = 0$. A composite
passes only when $a$ lies in its \emph{strong-liar set}, and Rabin's theorem
bounds that set at $(n-3)/4$, so a uniform witness gives
\begin{equation}
\label{eq:mr-fpr}
  \fprate \;\le\; \frac{1}{4}.
\end{equation}
A single round is therefore a $\Bernoulli{\BoolType}$ instance with $\fnrate =
0$ and $\fprate \le 1/4$.

To bring the $k$-round test under Axiom~1 (\Cref{ax:element-indep}), recast it
as a predicate on the witness pool: $\tilde{P}(a)$ is the verdict for witness
$a$, the pool $\{2, \ldots, n-2\}$ is the universe, and distinct witnesses are
distinct input elements. The $k$-round test is the conjunction $\tilde{P}(a_1)
\wedge \cdots \wedge \tilde{P}(a_k)$ over $k$ independently drawn witnesses,
exactly the operation of \Cref{sec:and-under-noise}.

\subsection*{The k-round bound}

Each round is one-sided with $\fnrate_i = 0$, so the conjunction has
$\fnrate_k = 0$: the test never calls a prime composite. All $k$ rounds test
the same $n$, so the only true input is $\bot$ (for a composite), and the
false-positive rate is the all-negative dual of the composition theorem,
\eqref{eq:composition-theorem}: the conjunction false-positives only if every
witness is a strong liar, so
\begin{equation}
\label{eq:mr-k-round}
  \fprate_k \;=\; \prod_{i=1}^{k} \fprate_i \;\le\; \left(\frac{1}{4}\right)^k.
\end{equation}
The equality is the product over the independently drawn witnesses (membership
of a uniform draw in the strong-liar set is independent across draws); the
inequality bounds each round by the Monier-Rabin value $\fprate_i \le 1/4$.
Both ingredients are needed: the size bound $|L_n| \le (n-3)/4$ from number
theory and the independent sampling that makes Axiom~1 apply. At $k = 5$ the
bound is below $10^{-3}$, and at $k = 20$ below $10^{-12}$, effectively
definitive.

\begin{algorithm}[ht]
  \caption{Miller-Rabin primality test ($k$ rounds)}
  \label{alg:miller-rabin}
  \SetKwInOut{Input}{input}
  \SetKwInOut{Output}{output}
  \SetKwProg{func}{function}{}{}
  \Input{$n$: odd integer $> 3$; \quad $k$: number of rounds}
  \Output{\textsc{composite} (definitive), or \textsc{probably-prime} (FPR $\le (1/4)^k$)}
  \func{\textsc{MillerRabin}($n$, $k$)}{
    write $n - 1 = 2^r d$ with $d$ odd\;
    \For{$i \gets 1$ \KwTo $k$}{
      choose $a$ uniformly from $\{2, \ldots, n-2\}$\;
      $x \gets a^d \bmod n$\;
      \lIf{$x = 1$ \textbf{or} $x = n-1$}{\textbf{continue}}
      \For{$j \gets 1$ \KwTo $r - 1$}{
        $x \gets x^2 \bmod n$\;
        \lIf{$x = n - 1$}{\textbf{break}}
      }
      \lIf{$x \neq n - 1$}{\Return \textsc{composite}}
    }
    \Return \textsc{probably-prime}\;
  }
\end{algorithm}

\subsection*{A worked composite: \texorpdfstring{$n = 341$}{n = 341}}

The number $341 = 11 \times 31$ is a Fermat pseudoprime to base $2$ ($2^{340}
\equiv 1 \pmod{341}$), so the Fermat test passes it, but Miller-Rabin does not.
With $n - 1 = 2^2 \cdot 85$, take $a = 2$: $2^{85} \equiv 32 \pmod{341}$ with
$32 \neq \pm 1$, and $32^2 \equiv 1 \pmod{341}$. A non-trivial square root of
$1$ has appeared, which a prime modulus forbids, so the witness $a = 2$
certifies $341$ composite. Miller-Rabin is stronger than Fermat precisely
because it inspects the square roots of $1$ along the way to $a^{n-1}$, not
just the endpoint.

\subsection*{The unification}

Miller-Rabin is not a separate technique. The single round is a
$\Bernoulli{\BoolType}$ instance with $\fnrate = 0$ and $\fprate \le 1/4$; the
$k$-round test is the conjunction of \Cref{sec:and-under-noise}; the bound
$\fprate_k \le (1/4)^k$ is \eqref{eq:mr-k-round}; and the design rule, choose $k
\ge \tfrac{1}{2}\log_2(1/\epsilon) = \log_4(1/\epsilon)$ to reach
false-positive target $\epsilon$, is that bound solved for $k$ ($k = 5$
reaches $10^{-3}$ and $k = 20$ reaches $10^{-12}$, the rounds quoted above). The number theory supplies the per-round $1/4$; the rest
is Boolean algebra over noisy bits. The broader theory of randomized
algorithms, including the Las Vegas versus Monte Carlo distinction, is
\Cref{ch:randomized-algorithms}; the companion one-sided example, the Bloom
filter, is \Cref{sec:bridge-to-sets}.
